\section{Recomendaciones y Conclusiones}

\subsection{Recomendaciones}

\subsubsection{Investigadores}

Quienes usen Astro-Denoise deben citar la versión específica de código y datos, respetar las licencias, y publicar sus propios parámetros y código. La reproducibilidad se construye con práctica colectiva.

\subsubsection{Comités de Evaluación}

Se requiere monitoreo anual de principios éticos, auditorías de acceso, y mecanismos claros para resolver conflictos de autoría. El marco ético debe actualizarse cada tres años para incorporar normativas emergentes.

\subsection{Riesgos y Mitigaciones}

Los tres riesgos identificados — irreproducibilidad, acceso inequitativo y gobernanza indefinida — ya tienen controles implementados. Riesgos futuros como cambio de plataforma, conflicto de intereses o degradación técnica se mitigan con licencia MIT (no privatizable), DOI permanente, forks comunitarios y revisiones técnicas anuales.

\subsection{Criterios de Éxito}

\begin{table}[H]
\centering
\small
\begin{tabular}{|p{3cm}|p{5cm}|}
\hline
\textbf{Criterio} & \textbf{Indicador} \\
\hline
Reproducibilidad & Artículos citan versión específica de código y datos \\
\hline
Equidad & Contribuidores de diversas regiones y tipos de institución \\
\hline
Acceso & Uso activo documentado en repositorios públicos \\
\hline
Gobernanza & Cero violaciones de licencia documentadas \\
\hline
\end{tabular}
\end{table}

\subsection{Conclusión}

El proyecto Astro-Denoise se alinea con la normativa colombiana (Ley 1581, Ley 2502, CONPES 4144, Minciencias), los estándares internacionales (FAIR, UNESCO) y los principios de ciencia responsable. Los riesgos identificados tienen controles implementados. Se recomienda mantener el compromiso de acceso abierto, la revisión ética periódica y la documentación de decisiones de diseño a lo largo del ciclo de vida del proyecto.

\bigskip

\textbf{Documento}: 2026-06-07 \hfill \textbf{Responsable}: Jheison Martinez Bolivar, Universidad La Salle
